\scanpage{6}
\uline{Proposition} On a donc des \deleted{\unsure{homo|les}} inclusions
\[
\mathcal E_0\unsure{\subset|=}\mathcal E_1\subset\mathcal E_2\subset\mathcal E_3\subset\cdots
\]
les $\mathcal E_i$ \unsure{étant|sont} des faisceaux loc. libres de rang $i$ pour $i\geq1$. D'ailleurs
\[
\mathcal E_0=\underline O_S
\]

\uline{Proposition} 2 $\mathcal E_0=\mathcal E_1=\underline O_S$, d'autre part pour $i\leq j$ l'injection $\mathcal E_i\longrightarrow\mathcal E_j$ est \deleted{loc$^{\rm t}$.} universellement injective. \inserted{Soit} \deleted{Le faisceau} $\mathcal L_n\simeq\mathcal E_{\deleted{\unsure{n+1}}n}/\mathcal E_{n-1}$ \inserted{(pour $n\geq1$, \deleted{\unsure{$\underline O_S$}}, donc $\mathcal L_i$ est inv si $i\ne1$), alors on a des isomorphismes} \deleted{est isomorphe à $\mathcal L^{\deleted{\unsure{\otimes}}}$, et $\mathcal E_{n+1}/\mathcal E_n$ est isomorphe à $\mathcal L^{\otimes n}$.} $\mathcal L_n\otimes\mathcal L_m\simeq\mathcal L_{m+n}$ si $n,m\left|\begin{smallmatrix}\geq0\\\ne1\end{smallmatrix}\right.$, \inserted{correspondant} \deleted{faisant de} à \inserted{une structure} sur $\coprod_{n\geq0}\mathcal L_n$ \deleted{\unsure{une}} d'Algèbre graduée comm.

Pour la première assertion, à cause de la platitude coh., on est ramené au cas d'un corps, où c'est trivial car on a une injection.

Notons que $\mathcal E=\varinjlim\mathcal E_n=f_*(\varinjlim\underline L^{\otimes n})$ est une $\underline O_S$-\uline{algèbre}, dont la connaissance d'ailleurs implique la connaissance de $\underline\Gamma_*^{\unsure{\mathcal E|\xi}}(f,\underline L)=\coprod\mathcal E_n$, car $\mathcal E_m\times\mathcal E_n\longrightarrow\mathcal E_{m+n}$ est précisément induit par la multiplication de $\mathcal E$. On a, en passant au gradué associé à la filtration croissante donnée, une algèbre $\coprod_{n\geq0}\underline{\mathcal L}_n$, où les $\underline{\mathcal L}_n\simeq\mathcal E_n/\mathcal E_{n-1}$ pour $n\geq0$, en modifiant pour $n=0$ par $\underline O_S$. Il faut prouver que $\underline{\mathcal L}_n\otimes\underline{\mathcal L}_m\longrightarrow\underline{\mathcal L}_{n+m}$, $n,m\left|\begin{smallmatrix}>0\\\ne1\end{smallmatrix}\right.$, \deleted{$n,m$} \deleted{\unsure{est}} bien isom., ou encore surjectif, \deleted{et pour} Si $n$ ou $m=0$ \deleted{\unsure{réciproquement si $n$ et $m$ sont entiers}} c'est trivial, sinon on a $n,m\geq2$, et pour voir.
% archive sheet number: 5
